\chapter*{Einleitung zur zweiten Lerneinheit}
\addcontentsline{toc}{chapter}{Einleitung zur zweiten Lerneinheit}

Nachdem die erste Lerneinheit vor allem Wiederholungs- bzw. Brückenkurscharakter hatte, geht es mit der vorliegenden zweiten Lerneinheit »richtig los«. Mit anderen Worten wird davon ausgegangen, dass die meisten Studierenden mit den hier präsentierten Inhalten noch nicht vertraut sind.

Es geht um eine Einführung in die Finanzmathematik. Aus mathematischer Sicht enthält diese Lerneinheit nicht viel Neues -- bei den hier vorgestellten Grundlagen der Zinseszins-, Renten- und Tilgungsrechnung kommen neben den Grundrechenarten noch Potenzen, Logarithmen und Wurzeln zum Einsatz, ferner muss man Gleichungen umstellen können usw. Neu sind dagegen die Begriffsbildungen wie Barwert, nachschüssige Rente oder Effektivzinssatz. Die Schwierigkeit bei der Anwendung der finanzmathematischen Verfahren besteht eher darin, die jeweils gegebene Anwendungssituation korrekt zu analysieren und zu entscheiden, welche Formeln anzuwenden sind. Am besten lernt man dies durch das eigenständige Lösen von Übungsaufgaben -- die angegebenen Lösungen sollte man erst nach der eigenen Bearbeitung ansehen.

Für an der Finanzmathematik besonders Interessierte soll noch angefügt werden, dass es in diesem Gebiet eigentlich erst nach den hier präsentierten Inhalten richtig »spannend« wird: Bei der mathematischen Behandlung von Aktien und Optionen müssen jedoch wahrscheinlichkeitstheoretische Methoden verwendet werden, so dass dies in der vorliegenden Lerneinheit nicht behandelt werden kann.
